\title{Parameter-Efficient Orthogonal Finetuning via Factorization}

\begin{abstract}
The widespread adoption of large foundation models comes with substantial computational demands, making their full retraining impractical. As such, finding effective methods for adapting these pre-trained models to novel tasks becomes critical. This study explores a systematic approach to finetuning—Orthogonal Finetuning (OFT)—to improve performance on downstream tasks. While OFT has shown strong generalization capabilities, its requirement for manipulating high-dimensional orthogonal matrices results in a considerable parameter footprint. To mitigate this, we first analyze OFT from the lens of information transfer, leading us to uncover several crucial attributes necessary for enhancing parameter efficiency. Motivated by the structure of the Cooley-Tukey fast Fourier transform, which streamlines information propagation, we introduce a compact orthogonal parameterization based on butterfly architectures. By integrating this scheme within OFT, we develop Orthogonal Butterfly~(BOFT), a new parameter-efficient finetuning method. BOFT generalizes OFT—treating it as a particular instance—and presents an expanded framework for orthogonal finetuning. We validate the efficacy of BOFT through extensive experiments involving large vision transformers, language models, and text-to-image diffusion architectures across diverse tasks in vision and language domains.
\end{abstract}

\section{Introduction}

State-of-the-art models such as ChatGPT~\cite{brown2020language,chen2021evaluating} and Stable Diffusion~\cite{rombach2022high} showcase the extraordinary generalization strengths inherent to large foundation models. This surge in model capabilities has been accompanied by an explosive growth in parameter counts

(\emph{e.g.}, GPT-3 is built with roughly 175B parameters), rendering the training of such models from the ground up increasingly unattainable for most researchers. Consequently, advancing the field hinges upon techniques for repurposing pre-trained foundation models rather than full retraining. The central challenge thus becomes how to swiftly and resourcefully adapt existing models for specific downstream applications. Broadly, there are three prominent strategies for efficient adaptation: \emph{model finetuning}~\cite{hulora2022,zaken2022bitfit,guo2020parameter,raffel2020exploring,qiu2023controlling,zhang2022adaptive,chavan2023one}, which leaves the original model architecture intact and updates only a selected parameter subset; \emph{adapter tuning}~\cite{rebuffi2017learning,houlsby2019parameter,pfeiffer2020adapterfusion,lin2020exploring,he2021towards}, which inserts new trainable parameters within the model structure; and \emph{prompt tuning}~\cite{li2021prefix,lester2021power}, which incorporates additional learned tokens at the input level. Among these, model finetuning stands out due to its simplicity, practical effectiveness, and its advantage of not incurring any added inference-time delays.

The core idea underpinning model finetuning is to maintain similarity between the pretrained and finetuned models according to particular metrics, thereby preserving the knowledge acquired during pretraining. For example, prevalent approaches in model finetuning typically employ a low learning rate, an ad hoc method that implicitly minimizes the deviation between the initial and updated model weights. Considering that weight matrices are structured, a more systematic strategy involves conserving the relational properties of these matrices, specifically the pairwise angles among neurons. This perspective gives rise to Orthogonal Finetuning~(OFT)~\cite{qiu2023controlling}, an innovative framework in which all neurons within the same layer are collectively transformed via a single orthogonal matrix. This guarantees the theoretical preservation of pairwise neuronal angles throughout the finetuning phase. Although OFT has shown notable improvements in both generalization and convergence when applied to the finetuning of text-to-image diffusion models~\cite{qiu2023controlling}, a significant drawback lies in the sheer volume of trainable parameters required, stemming from the large dimensionality of the orthogonal matrices involved. In response, OFT utilizes a block-diagonal formulation to curtail parameter count. Nevertheless, this gain in parameter efficiency introduces certain limitations—chief among them being the imposition of a rigid sparsity pattern and the independent application of orthogonal transformations across distinct blocks. While this block-diagonal sparsity pattern achieves parameter savings, it may simultaneously generate suboptimal inductive biases (\emph{e.g.}, diminishing the expressive capacity and preventing the matrix from representing conventional linear operations), and the optimal selection of such sparsity patterns remains unresolved.

The central challenge lies in constructing a parameter-efficient dense orthogonal matrix. At first, this appears intractable, as forming a dense orthogonal matrix in $d$ dimensions traditionally requires $\mathcal{O}(d^2)$ parameters. However, our work introduces a novel solution: synthesizing a dense orthogonal matrix through the composition of several factorized sparse matrices. This methodology builds on the observation that the parameter count may be curtailed by exchanging additional computational time for reduced spatial requirements. By encoding the orthogonal matrix as a product of sparse matrices, it becomes necessary to repeatedly compute these products during each training step. 

To conceptualize our approach to matrix factorization, we reframe the process of creating a dense orthogonal matrix as an information propagation task. Within this perspective, the construction via sparse matrix multiplications can be analogized to information transfer across a grid-like graph structure. Viewing the problem through this lens enables a range of sparse matrix factorization strategies that limit trainable parameters, all while retaining sufficient capacity to model dense matrices.

To further advance parameter efficiency within our information transmission schema, we are inspired by the butterfly patterns of the Cooley-Tukey fast Fourier transform algorithm~\cite{cooley1965algorithm}, where information is exchanged effectively between nodes~\cite{klasing1994broadcasting}. The butterfly network implicit in the Cooley-Tukey procedure yields what is termed butterfly factorization—a particular kind of sparse matrix decomposition. Utilizing butterfly factorization, we can express a $d$-dimensional dense matrix as a sequence of $\mathcal{O}(\log d)$ sparse factors, each having only $\mathcal{O}(d)$ nonzero components. Ensuring that every sparse matrix factor remains orthogonal, we achieve a dramatically more efficient orthogonal parameterization: this representation contains only $\mathcal{O}(d\log d)$ parameters, a substantial improvement over the conventional approach. Building on this orthogonal parameterization, we introduce a new, parameter-efficient finetuning paradigm called Orthogonal Butterfly~(BOFT). Notably, BOFT generalizes and includes OFT as a particular instance, establishing a broad framework for orthogonal finetuning. 

A key point of similarity between the block-diagonal and butterfly matrix structures is their shared sparsity, which both exploit to reduce the count of tunable parameters within orthogonal matrices. In contrast to LoRA~\cite{hulora2022}, which adopts low-rank matrix structures, our method utilizes sparse matrices; both approaches represent core families of structured matrices~\cite{candes2011robust} aimed at promoting parameter efficiency.

In contrast to the block-diagonal configuration employed by OFT, which seeks a balance between model expressiveness and structural regularity, BOFT adopts the butterfly structure to provide a smoother continuum between matrices representing the full orthogonal group (maximizing expressivity) and the identity (emphasizing regularity). This design makes it possible to restrict the solution space to a narrower, more tractable hypothesis class within the orthogonal group, potentially benefitting downstream applications. Given the butterfly structure's centrality in numerous efficient linear transformations—such as those underlying the discrete Fourier and cosine transforms—our structured methodology for parameter-efficient modeling is expected to impart a natural inductive bias. Such bias could enhance the model's ability to generalize while curbing overfitting. The central premise is that, while the butterfly structure can replicate many standard linear transforms, this capability is absent in the block-diagonal approach of OFT. Our specific contributions include:

\begin{itemize}[leftmargin=*,nosep]
    \item We analyze parameter efficiency in orthogonal finetuning by proposing a novel information transmission framework. This formulation reconceptualizes the challenge of building a parameter-efficient dense orthogonal matrix as an information flow problem over a grid-structured graph.
    \item Drawing on the butterfly structure at the heart of the Cooley-Tukey algorithm, we introduce Orthogonal Butterfly—a parameter-efficient orthogonal finetuning strategy framed within the information transmission perspective.
    \item We offer theoretical justification for BOFT’s effectiveness, demonstrating how it achieves substantial reductions in trainable parameters without compromising expressive capacity or training stability. Furthermore, the matrix factorization inherent to BOFT naturally supports interesting forms of weight interpolation (see Figure~\ref{fig:boft_interpolation}).
    \item For the first time, we extend the application of orthogonal finetuning to a broad suite of tasks beyond controllable text-to-image generation~\cite{qiu2023controlling}, highlighting its promise as a universal finetuning technique. In particular, BOFT is deployed across diverse adaptation settings—from computer vision to natural language processing—and achieves superior performance over the current state-of-the-art, illustrating its remarkable parameter-efficiency and capacity for generalization.
\end{itemize}

\textbf{Parameter-efficient finetuning (PEFT)}. As foundation models (\emph{e.g.}, \cite{brown2020language,radford2021learning,rombach2022high,kirillov2023segment}) grow in scale and capacity, developing methods for adapting them to downstream tasks while minimizing the number of tunable parameters has become a focal area of research~\cite{treviso2023efficient,ding2023parameter,lialin2023scaling}. Among the diverse approaches within PEFT~\cite{houlsby2019parameter,aghajanyan2020intrinsic,hulora2022,edalati2022krona,wang2022adamix,gheini2021cross,zaken2022bitfit,guo2020parameter,sung2021training,ansell2022composable,lester2021power,li2021prefix,vu2022spot,he2021towards,mao2021unipelt,karimi2021compacter,liu2022few,sung2022lst,chen2022parameter,jia2022visual,chen2022adaptformer,zhang2022neural,jie2023fact,lian2022scaling,luo2023towards}, methods relying on reparameterization~\cite{aghajanyan2020intrinsic,hulora2022,edalati2022krona} are most pertinent to the present study. A prominent example is LoRA~\cite{hulora2022}, which finetunes pretrained weights by incorporating a product of two low-rank matrices, thereby attaining competitive outcomes on various natural language tasks. As LoRA sets a fixed rank for all low-rank matrices, subsequent research~\cite{zhang2022adaptive,valipour2022dylora,zhang2023increlora} has investigated adaptive rank selection across layers to optimally utilize parameter resources. Owing to the straightforwardness of low-rank reparameterization, this paradigm has achieved widespread adoption~\cite{dettmers2023qlora,zi2023delta,chavan2023one}. Motivated by hyperspherical energy analyses of generalization~\cite{liu2018learning,liu2021orthogonal}, \cite{qiu2023controlling} introduce orthogonal finetuning (OFT), which offers an alternative strategy for refining text-to-image diffusion models. OFT optimizes an orthogonal transformation within each layer, leading to improved generalization properties and more consistent training behavior compared to LoRA. Nevertheless, OFT typically entails a larger set of trainable parameters relative to LoRA, highlighting the value of making OFT more parameter-efficient. Furthermore, the suitability of OFT for a broader variety of adaptation challenges beyond text-to-image diffusion~\cite{qiu2023controlling} remains an open question. BOFT seeks to advance OFT's parameter efficiency using butterfly factorization, thereby facilitating a wider exploration of orthogonal finetuning across general adaptation scenarios.

\textbf{Butterfly structures}. The radix-2 Cooley-Tukey algorithm transforms an $N$-point discrete Fourier transform into two smaller $\frac{N}{2}$-point transforms through a recursive procedure, naturally yielding a butterfly pattern—mathematically represented as a product of several sparse matrices known as a butterfly matrix. Butterfly matrices have been employed as parameterizations for orthogonal matrices, eliminating the need for pivoting in Gaussian elimination and boosting computational efficiency~\cite{parker1995random}. They have also found application in stabilizing recurrent neural network training~\cite{jing2017tunable} and in kernel approximation methods~\cite{munkhoeva2018quadrature}. Research by \cite{dao2019learning,dao2019kaleidoscope} has leveraged butterfly parameterizations to design fast linear transformations, while \cite{chen2022pixelated,dao2022monarch} have used butterfly matrices to improve neural network training efficiency. Their utility extends to fast matrix-vector computations~\cite{michielssen1996multilevel,o2010algorithm}, approximating data-sparse matrices~\cite{li2015butterfly}, and efficient information transmission in network contexts~\cite{klasing1994broadcasting,li2003linear,rajasingh2016transmission}. Distinct from earlier studies, our approach centers on leveraging butterfly structures specifically to increase the parameter efficiency of OFT.

\section{An Information Transmission View on Orthogonal Finetuning}

We begin by reviewing foundational aspects of OFT. In order to adapt a pretrained weight matrix, OFT reformulates the updated weights as the product of a trainable orthogonal matrix and the static pretrained weights. Unlike LoRA, which updates weights by adding a low-rank matrix, OFT relies on a multiplicative update involving an orthogonal matrix. LoRA’s parameter-efficiency is attributed to its use of low-rank structures, whereas OFT achieves similar efficiency by employing a block-diagonal orthogonal configuration~\cite{qiu2023controlling}; specifically, smaller diagonal block sizes correspond to greater parameter-efficiency in OFT. Figure~\ref{fig:comp_lora} illustrates this comparison. The underlying rationale for utilizing orthogonal transformations in weight finetuning is to maintain the angular relationships between neuron pairs~\cite{liu2018learning,liu2021orthogonal,qiu2023controlling}, thereby conserving much of the semantic information acquired during pretraining. Formally, OFT learns an orthogonal matrix $\bm{R}\in\mathbb{R}^{d\times d}$ for a pretrained linear transformation $\bm{W}^0\in\mathbb{R}^{d\times n}$, modifying the computation from $\bm{z}=(\bm{W}^0)^\top\bm{x}$ to $\bm{z}=(\bm{R}\bm{W}^0)^\top\bm{x}$, where $\bm{x}\in\mathbb{R}^{d}$ is the input vector and $\bm{z}\in\mathbb{R}^{n}$ denotes the output. To realize zero initialization, the matrix $\bm{R}$ is initialized as the identity matrix. Maintaining the orthogonality of $\bm{R}$ during the finetuning process, we adopt the Cayley transform~\cite{qiu2023controlling,liu2021orthogonal}, parameterizing $\bm{R}$ as $\bm{R}=(\bm{I}+\bm{Q})(\bm{I}-\bm{Q})^{-1}$, where $\bm{Q}$ is a skew-symmetric matrix satisfying $\bm{Q}=-\bm{Q}^\top$. For enhanced parameter-efficiency, a block-diagonal form is used by specifying $\bm{R}$ as $\text{diag}(\bm{R}_1,\bm{R}_2,\cdots,\bm{R}_r)$, in which each $\bm{R}_i\in\mathbb{R}^{b\times b}$ represents a small orthogonal block and $br=d$. However, this block-diagonal design introduces an inherent constraint: each neuron's dimensions (i.e., columns of $\bm{W}^0$) are partitioned into $r$ separate groups, and each group is transformed via a different orthogonal block, independent of the neuron's functional meaning. Although this method has proven empirically effective, the rationale for grouping neuron dimensions solely by index is weak, thereby motivating the pursuit of dense orthogonal matrices. Consequently, a central question emerges: \emph{Is it feasible to devise a dense orthogonal matrix while retaining parameter-efficiency?}

To tackle this issue, we suggest representing a dense orthogonal matrix as the product of several sparse orthogonal matrices. To systematically and intuitively analyze the sparsity structures emerging from orthogonal matrix factorizations, we conceptualize the task of constructing a dense orthogonal matrix in OFT as analogous to an information transmission scenario. Specifically, generating a dense matrix $\bm{R}\in\mathbb{R}^{d\times d}$ as a product of $m$ square matrices, i.e., $\thickmuskip=2mu \medmuskip=2mu \bm{R}=\bm{B}_m\bm{B}_{m-1}\cdots\bm{B}_1$, can be seen as relaying information through a $d\times (m+1)$ node grid, depicted in Figure~\ref{fig:info_trans}. This information transmission metaphor stems from interpreting a $d$-dimensional dense square matrix as encoding all-to-all connectivity between two sets of $d$ nodes. For $\bm{R}$, a zero entry $R_{ij}$ signifies a missing information pathway from node $j$ to node $i$, whereas a nonzero value allows for information flow. Consequently, expressing the dense matrix $\bm{R}$ as the sequence $\bm{B}_m\bm{B}_{m-1}\cdots\bm{B}_1$ can be viewed as a series of information exchanges determined by the graph structures generated by each $\bm{B}_i$. Here, information first traverses the connections dictated by $\bm{B}_1$ and finally passes through $\bm{B}_n$. To illustrate, Figure~\ref{fig:info_trans} showcases the decomposition $\bm{R}=\bm{B}_5\bm{B}_4\bm{B}_3\bm{B}_2\bm{B}_1$ with their sparsity layouts and associated graph. This graph is the result of “unrolling” the sequence of matrix multiplications. Each matrix $\bm{B}_i$ functions as a connectivity map from nodes at level $i$ to nodes at level $i+1$. More precisely, the $(j_1, j_2)$-th entry of $\bm{B}_i$ indicates the presence (nonzero) or absence (zero) of a directed edge from the $j_2$-th node of the $i$-th level to the $j_1$-th node in the $(i+1)$-th level. For $\bm{B}_5\bm{B}_4\bm{B}_3\bm{B}_2\bm{B}_1$ to yield a dense $\bm{R}$, every node at the initial level must be able to reach all nodes at the sixth level via information flow. If we restrict attention to $\bm{R}=\bm{B}_4\bm{B}_3\bm{B}_2\bm{B}_1$, which connects the nodes at the first and fifth levels, we observe that node $1$ does not have a path to node $3$, indicating that the sparsity pattern in $\bm{B}_4\bm{B}_3\bm{B}_2\bm{B}_1$ must contain a zero at position $(3,1)$. Similar observations hold for factorizations such as $\bm{R}=\bm{B}_3\bm{B}_2\bm{B}_1$ and $\bm{R}=\bm{B}_2\bm{B}_1$, where the graphs and sparsity patterns remain consistent. More generally, to ensure that a matrix $\bm{R}\in\mathbb{R}^{d\times d}$ is dense, the collection of factor matrices $\bm{B}_m,\ldots,\bm{B}_1$ must define directed edges on a $d\times(m+1)$ grid. In this structure, each directed edge only links two nodes from consecutive levels (columns), and information originating at any node in the first level must be capable of reaching any node at the $(m+1)$-th level.

Viewing the problem through the lens of information transmission offers deeper insight into the sparsity pattern exhibited by the factorization matrices in OFT. Figure~\ref{fig:blk_diag} illustrates the block-diagonal arrangement of $\bm{R}$ that is characteristic of conventional OFT. Although such a structure leads to a decrease in the quantity of trainable parameters, it also prevents $\bm{R}$ from becoming a densely populated matrix. Our objective is to assemble a dense orthogonal matrix by composing $m$ sparse orthogonal matrices, striving to minimize the number of effective trainable parameters involved. Within the information transmission perspective, our general requirements can be described as: \emph{(i)} \textbf{dense connectivity}, whereby each node at the input stage must maintain at least one connection to every node at the output stage; and \emph{(ii)} \textbf{minimum free edges}, seeking to keep the sum of edges, given the orthogonality requirement, as small as possible. The orthogonality condition imposes nuanced restrictions on the admissible edges between consecutive layers. For instance, to guarantee each matrix $\bm{B}_i$ has full rank—a prerequisite for orthogonality—it is necessary to establish $d$ edges that map all nodes at the $i$-th stage to all nodes at the $(i+1)$-th stage, thereby setting the minimum number of inter-level edges at $d$ (see, for example, Figure~\ref{fig:info_trans} where $d=4$). Since these $d$ edges are dictated by the requirement of orthogonality and do not involve adjustable parameters (their values are fixed, such as $\pm1$ in a $d\times d$ orthogonal matrix with only $d$ non-zero elements), they should be excluded from the trainable edge count. The orthogonality constraint, by reducing the required number of parameters compared to unconstrained full matrices, also introduces further dependencies among the edge configuration. For example, in the case of a $2\times 2$ orthogonal matrix, if the entry at position $(1,1)$ is zero, orthogonality forces the entry at position $(2,2)$ to be zero as well (\emph{i.e.}, omitting one edge can necessitate omitting another). Consequently, orthogonality may sometimes necessitate adding or eliminating specific connections within any admissible edge set. Examining the pattern of non-zero elements in the constructed orthogonal matrix, the information transmission perspective becomes particularly valuable in OFT contexts, since the only relevant components of $\bm{R}$ are its non-zero, trainable entries, while their actual numerical values are inconsequential.

Establishing a fully dense connection between two levels in a network involves $\mathcal{O}(d^2)$ edges (corresponding to a single dense orthogonal matrix), resulting in $d^2-d$ total edges (for example, when $d=4$, this amounts to 12 edges). Figure~\ref{fig:info_trans} presents a matrix factorization that requires only 10 edges in total, demonstrating an improvement in edge efficiency over a dense orthogonal matrix. This approach allows the exploration of OFT's parameter-efficiency from a graphical viewpoint, and the proposed framework facilitates the systematic construction of valid matrix factorizations. Motivated by a structure from the Cooley-Tukey algorithm known as butterfly graphs~\cite{cooley1965algorithm}, which effectively links $d$ sources and $d$ destinations using just $\mathcal{O}(d\log d)$ edges, we note that while the configuration in Figure~\ref{fig:info_trans} needs 10 edges for complete connectivity, a butterfly network achieves this with only 8 edges. In what follows, we discuss how employing the butterfly topology can further optimize parameter efficiency.

\section{Orthogonal Parameterization by Butterfly Factorization}

Originally introduced within the Cooley-Tukey algorithm to accelerate the computation of the fast Fourier transform, the butterfly structure plays a key role in facilitating efficient global information propagation—whereby local frequency modifications yield global effects in the spatial domain. This concept resonates with our information transmission objective: ensuring all first-level nodes can communicate with any node at the output layer. The butterfly network is also renowned as a computer network topology~\cite{solihin2015fundamentals,li2011linear}, recognized for its efficacy in data exchange. For $k\geq 2$, where $k$ is a power of $2$, we define the \emph{butterfly factor} $\bm{B}^F(k)$ as

\begin{equation}\label{eq:but_factor}
\begin{aligned}
    \bm{B}^F(k)=\begin{bmatrix}
    \text{diag}(\bm{d}_1) & \text{diag}(\bm{d}_2) \\
    \text{diag}(\bm{d}_3) & \text{diag}(\bm{d}_4)
    \end{bmatrix}\in\mathbb{R}^{k\times k},
\end{aligned}
\end{equation}

where each $\bm{d}_i \in \mathbb{R}^{\frac{k}{2}}$ for all $i$. Setting $d=2^N$, the $d$-dimensional \emph{butterfly matrix} $\bm{B}(d)\in\mathbb{R}^{d\times d}$ is defined recursively as follows:

\begin{equation}
\begin{aligned}
    \!\!\bm{B}(d)=\tilde{\bm{B}}(d,d)\!\cdot\!\begin{bmatrix}
    \bm{B}_1(\frac{d}{2})\!\!\!\!\! & \bm{0} \\
    \bm{0}\!\!\!\!\! &\bm{B}_2(\frac{d}{2})
    \end{bmatrix}= \tilde{\bm{B}}(d,d)\tilde{\bm{B}}(d,\frac{d}{2})\cdots\tilde{\bm{B}}(d,2),
\end{aligned}
\end{equation}

where $\bm{B}_1(\frac{d}{2})$ and $\bm{B}_2(\frac{d}{2})$ denote butterfly matrices each of dimension $\frac{d}{2}$. We introduce the \emph{butterfly component} as $\thickmuskip=2mu \medmuskip=2mu \tilde{\bm{B}}(d,k)=\text{diag}(\bm{B}^F_1(k),\cdots,\bm{B}^F_{d/k}(k))$, representing a $d\times d$ block-diagonal matrix in which every diagonal block has dimension $k$. Here, each $\bm{B}^F_i(k)$ is a butterfly factor as outlined in Equation~\ref{eq:but_factor}. This setup enables the construction of an orthogonal matrix through butterfly matrix parameterization. To do so, it suffices to make sure each multiplicative element $\tilde{\bm{B}}(d,k)$ within the butterfly decomposition $\bm{B}(d)$ remains orthogonal. 

Focusing initially on $\tilde{\bm{B}}(d,2)$, a block-diagonal matrix with $2\times 2$ blocks, orthogonality is straightforwardly achieved by employing the Cayley transform, or equivalently a $2$-dimensional rotation, for every block, following the strategy in \cite{liu2021orthogonal,qiu2023controlling}. Since each butterfly component’s nonzero structure is a rearrangement of that of $\tilde{\bm{B}}(d,2)$, all butterfly components are readily parameterized to be orthogonal matrices. This technique leads to an efficient representation of orthogonal matrices that consists of numerous $2\times 2$ orthogonal blocks.

Following the generalization proposed in \cite{chen2022pixelated}, we extend the concept by introducing the block butterfly component $\tilde{\bm{B}}^b(d,k)$, wherein every entry of each $\bm{d}_i$ is replaced by a $b\times b$ matrix. To ensure orthogonality of the block butterfly component $\tilde{\bm{B}}^b(d,2)$, every $2b\times 2b$ block is individually parameterized as an orthogonal matrix. The nonzero configuration for other components, i.e., $\tilde{\bm{B}}^b(d,k)$ with $k>2$, is a block-wise permutation of the structure in $\tilde{\bm{B}}^b(d,2)$, thus allowing them to be similarly parameterized as orthogonal matrices. Aggregating these constructions, the forward propagation in BOFT can be expressed as

\begin{equation}
    \begin{aligned}\nonumber
    \bm{z}=\big(\bm{R}(m,b)\cdot\bm{W}^0\big)^\top\bm{x},~~~~\text{s.t.}~\bigg{\{}\bm{R}(m,b)=\prod_{i=1}^m \tilde{\bm{B}}^b_{(d,i)}~~\&~~\underbrace{\big(\tilde{\bm{B}}^b_{(d,j)}\big)^\top\tilde{\bm{B}}^b_{(d,j)}=\tilde{\bm{B}}^b_{(d,j)}\big(\tilde{\bm{B}}^b_{(d,j)}\big)^\top\!=\bm{I}_d}_{\forall j\in [1, m]}\bigg{\}},
    \end{aligned}
\end{equation}

Here, for notational clarity, we use $\tilde{\bm{B}}^b_{(d,i)}$ as an abbreviation for $\tilde{\bm{B}}^b(d,2^{m - i+1})$, with $\bm{I}_d$ representing the $d$-dimensional identity matrix. The matrix $\bm{R}(m,b)\in\mathbb{R}^{d\times d}$ is a product of $m$ orthogonal butterfly blocks. For ease of reference, we label BOFT that utilizes $\bm{R}(m,\frac{b}{2})$ as BOFT($m$,$b$), where $b \geq 2$. In the special case where $\thickmuskip=2mu \medmuskip=2mu m=1$, BOFT($1$,$b$) simplifies to the block-diagonal OFT~\cite{qiu2023controlling} with blocks of size $b$. If the block size equals the full dimension, i.e., BOFT($1$,$d$), we recover the original OFT~\cite{qiu2023controlling} consisting of an unconstrained orthogonal matrix. Using a block butterfly matrix, such that BOFT($\log \frac{2d}{b}$,$b$) employs $\bm{B}^{\frac{b}{2}}(d)$ as $\bm{R}$, the construction produces a fully dense orthogonal matrix. Generally, BOFT($m$,$b$) incorporates $\frac{1}{2}(b-1)dm$ effective trainable parameters for adapting a $d\times n$ linear layer. If instead the classic butterfly structure is used, i.e., $m = \log d, b = 2$, the parameter count scales as $\mathcal{O}(d\log d)$. This is much lower than the $\mathcal{O}(d^2)$ parameter complexity for a fully dense OFT, and the $\mathcal{O}(bd)$ for a block-diagonal OFT with $r$ blocks. Notably, achieving a dense orthogonal matrix with original OFT requires setting $b = d$, whereas BOFT can realize a dense orthogonal matrix using any suitable $b$ value.

\textbf{BOFT Identity Initialization}. Fine-tuning typically begins from a precise copy of the pretrained parameters to ensure minimal drift from the original model. For instance, LoRA initializes its low-rank adaptations to zero. In the context of BOFT, all butterfly layers are initially set as identity matrices (i.e., skew-symmetric matrices used in the Cayley transformation are all initialized to zero).

\textbf{Multiplicative Dropout in BOFT}. LoRA~\cite{hulora2022} includes a Dropout operation for the low-rank updates to reduce overfitting, and standard Dropout~\cite{srivastava2014dropout} functions well in that setting. However, BOFT’s multiplicative update mechanism precludes using conventional Dropout directly. To address this, we introduce a new multiplicative Dropout for BOFT: since $\bm{R}(m,b)$ is comprised of $m$ butterfly matrices, which themselves can be rearranged into $2b\times 2b$ block-diagonal orthogonal matrices, we first randomly select a proportion $p_1$ of the butterfly blocks, and then randomly select $p_2$ of the diagonal blocks within each chosen butterfly. These selected diagonal blocks are then replaced by identity matrices.

\section{Intriguing Insights and Discussions}

\textbf{Expressivity of BOFT}. While the butterfly architecture supplemented with permutations is capable of exactly modeling several well-known fast linear transforms~\cite{dao2019learning,dao2019kaleidoscope} (such as the fast Fourier and Hadamard transforms), the extent to which our orthogonal butterfly matrix can approximate an arbitrary orthogonal matrix is yet to be fully explored. To investigate this, we carry out a simulation aimed at approximating a randomly sampled dense orthogonal matrix~\cite{anderson1987generation} of dimensions $512\times 512$. Figure~\ref{fig:para_eff} reports outcomes averaged over 10 independent random initializations. The vertical axis measures approximation error, while the horizontal axis corresponds to the number of effective trainable parameters. Each colored curve illustrates the performance of BOFT for a given block size; the leftmost point along each curve represents the error attained by BOFT($1$,$b$), i.e., the standard block-diagonal OFT with block size $b$. Empirically, BOFT consistently demonstrates superior parameter efficiency over OFT. For instance, BOFT($9$,$2$) achieves higher expressiveness than BOFT($1$,$16$), while requiring significantly fewer parameters. In general, reducing $b$ and increasing $m$ enhances the parameter efficiency of BOFT; as an example, BOFT($6$,$4$) attains an approximation error similar to BOFT($2$,$16$) but with far fewer parameters. These findings reflect the fact that the butterfly matrix forms a more constrained, structured subset within the set of orthogonal matrices relative to a block-diagonal construction, thereby ensuring that BOFT has provably greater expressiveness than an OFT with equal block size.

\begin{theorem}[Expressivity of BOFT]\label{thm:exp_boft}
    BOFT possesses greater expressivity compared to OFT when block sizes are matched. To approximate every orthogonal matrix of dimension $d$ using butterfly matrices, one may use products such as $\bm{B}_{d-1,1}(d)\bm{B}^\top_{d-1,2}(d)\cdots\bm{B}_{1,1}(d)\bm{B}^\top_{1,2}(d)$, where every $\bm{B}_{i,j}(d)$ is itself a butterfly matrix.
\end{theorem}

Based on Theorem~\ref{thm:exp_boft}, BOFT can be straightforwardly extended so that the final orthogonal matrix is generalized to $\thickmuskip=2mu \medmuskip=2mu \bm{R}^G(m_1,b_1,m_2,b_2,l)=\bm{R}_{l,1}(m_1,b_1)\bm{R}_{l,2}^T(m_2,b_2)\cdots\bm{R}_{1,1}(m_1,b_1)\bm{R}_{1,2}^T(m_2,b_2)$, where $\bm{R}_{i,j}^T(m,b)$ refers to the orthogonal matrices defined in BOFT. Particularly, when $m_1=m_2=\log d$, $b_1=b_2=2$, and $l=d-1$, the matrix $\bm{R}^G(m_1,b_1,m_2,b_2,l)$ is expressive enough to cover the entire orthogonal group—a construction referred to as the kaleidoscope hierarchy~\cite{dao2019kaleidoscope}. Nevertheless, higher expressivity does not guarantee improved performance during finetuning. In practice, full finetuning, while universally expressive, often yields suboptimal outcomes. Therefore, the fundamental challenge in model finetuning involves balancing expressiveness and regularization. By enriching the parameter space for finetuning with intrinsic structural priors, BOFT provides a mechanism to negotiate a more optimal compromise between model expressivity and regularity.

\textbf{Spectral properties}. Compared to LoRA, orthogonal finetuning tends to achieve superior spectral characteristics due to its exact preservation of the spectral norm associated with the pretrained weight matrix $\bm{W}^0$. This preservation becomes evident through the lens of singular value decomposition, expressed as $\bm{W}^0 = \bm{U}\bm{\Sigma}\bm{V}^\top$, with $\bm{U},\bm{V}$ representing orthogonal matrices and $\bm{\Sigma}$ being the diagonal matrix of singular values. In both OFT and BOFT, an orthogonal matrix $\bm{R}$ is applied on the left, resulting in finetuned weights of the form $\bm{R}\bm{U}\bm{\Sigma}\bm{V}^\top$. Importantly, this operation leaves the maximum singular value (i.e., the spectral norm of $\bm{W}^0$) unchanged. Such maintenance of the spectral norm has been shown to significantly promote training stability and improved generalization~\cite{miyato2018spectral, yoshida2017spectral}. Additional relevant mathematical properties are detailed in Appendix~\ref{app:sn_property}.

\textbf{Orthogonal finetuning as learning bilinear similarity}. BOFT may also be interpreted through the framework of bilinear similarity, specifically as $\bm{w}^0_i\bm{R}\bm{x}$, where $\bm{w}^0_i$ designates the $i$-th neuron (i.e., the column vector) in $\bm{W}^0$. Within this perspective, BOFT corresponds to training a bilinear similarity matrix $\bm{R}$ under a strict orthogonality constraint, fundamentally aligning this approach with concepts in distance metric learning~\cite{xing2002distance} and bilinear forms~\cite{roman2005advanced}.

\textbf{Inductive bias and generalization in BOFT}. The structure of $\bm{R}(m,b)$ in BOFT typically forms a particular subset within the orthogonal group, thereby constraining the hypothesis space and introducing an inductive bias. We propose that the specific structural bias conferred by butterfly factorization is advantageous for generalization because it mirrors shared patterns found in established linear transformations~\cite{dao2019kaleidoscope}, including the discrete Fourier transform, discrete sine/cosine transform, and the Hadamard transform. Furthermore, the inherent sparsity of the matrix factorization in BOFT may contribute additional implicit inductive biases~\cite{gunasekar2017implicit, li2018algorithmic, liu2019neural}.

\textbf{Comparison to butterfly-based sparse training}. Numerous studies~\cite{dao2019learning,dao2019kaleidoscope,chen2022pixelated,dao2022monarch} have explored sparse training approaches leveraging butterfly parameterizations. These prior works generally involve direct reparameterization of neural network weight matrices using butterfly structures and conducting training from scratch. \cite{dao2022monarch}, in particular, examines the process of finetuning by initially projecting pretrained weights onto a modified butterfly matrix and subsequently refining these components for specific downstream objectives. In contrast, BOFT introduces a novel finetuning approach, where layer-wise shared weight matrices are used to transform the weights.

\input{rebuttal-results}

\section{Concluding Remarks and Limitations}

This paper introduces Orthogonal Butterfly, a broadly applicable and parameter-efficient finetuning methodology for foundation models that builds upon the butterfly matrix structure. Central to our approach is the observation that enhanced parameter efficiency can be attained by representing a dense orthogonal matrix as a product of several sparse orthogonal matrices. To facilitate the search for suitable matrix decompositions, we design a graphical framework based on information transmission principles. Within this framework, we identify the butterfly structure as a compelling solution for sparse orthogonal matrix factorization. Our empirical analyses establish that BOFT is effective in finetuning large-scale language, vision, and text-to-image generative models, further showcasing its advantages as a versatile model adaptation method.

Nonetheless, BOFT is not without drawbacks. Specifically, the computational overhead during training is somewhat higher compared to OFT, as the ultimate orthogonal transformation in BOFT results from multiplying several orthogonal matrices. Addressing ways to decrease BOFT’s training cost remains an unresolved challenge. It is noteworthy, however, that once finetuning concludes, the orthogonal matrices learned by BOFT may be fused with the original model weights, thereby incurring no additional inference overhead. Furthermore, it is still unclear if the butterfly network structure offers optimal information transmission efficiency. Our graphical information transmission perspective invites connections to research in computer networking, where topological efficiency for data transmission is rigorously analyzed. We anticipate that the exploration of alternative, potentially more effective network topologies may inform future construction of dense orthogonal matrices.

\end{document}